\subsection{Recent History and Observation Heterogeneity Defined the Boundary}

\begin{figure*}[t]
\centering
\includegraphics[width=\textwidth]{figures/Fig5_final.pdf}
\caption{Observation heterogeneity. Route-level evidence and observation
classes show robust survivors, fragile boundaries, and stable failures across
the 19 observations. The metadata association panel is descriptive and
sensitivity-audited, but it is not a causal recording-condition explanation.}
\label{fig:heterogeneity}
\end{figure*}

Recent state-path history produced an intermediate result. Fourteen of the 19
observations exceeded the shuffled-history median, but only 6 of the 19
exceeded the q95 null threshold, and the sign and ordering of the effect were
not stable across observations. Thus, recent history was informative in an
observation-specific way, but it did not define a universal memory rule.

The observation-level synthesis showed that heterogeneity was reproducible and
observation-specific rather than erased by pooling. Thirteen observations were
classified as robust survivor/diffuse-positive cases, two as narrow-setting
rescue boundaries, two as stable non-survivors, one as a one-scale survivor,
and one as a robust survivor without the diffuse gate
(Fig.~\ref{fig:heterogeneity}). This heterogeneity constrains the manuscript
claim: T1 is common in this dataset, but the support, form, and explanatory
status of the residual vary across observations.

\begin{table*}[t]
\caption{Evidence-to-claim summary. Each claim is bounded by the tested
definition and should not be upgraded to the stronger interpretation listed in
the final column.}
\label{tab:evidence_claims}
\centering
\begingroup
\footnotesize
\setlength{\tabcolsep}{4pt}
\renewcommand{\arraystretch}{1.15}
\begin{tabular}{@{}p{0.18\textwidth}p{0.19\textwidth}p{0.27\textwidth}p{0.28\textwidth}@{}}
\toprule
Claim & Evidence & Allowed interpretation & Unsupported stronger claim \\
\midrule
Local non-affine survival & 14/19 both-scale and 15/19 any-scale survival; full-pipeline pseudo-event calibration $p_{\mathrm{both}}=0.000999$ at $B=1000$ & T1 is common after local affine subtraction & T1 is universal, causal, or a completed mechanism \\
Detrending challenge & Centered 14/19, past-only 11/19, no rolling detrending 13/19 & T1 survival is not erased by the challenge & The exact 14/19 count is preprocessing-invariant \\
Affine-fit conditioning QC & 38/38 observation-scale combinations passed; median condition number 2.37 & Local affine subtraction is numerically defensible in the sampled QC & All preprocessing artifacts are ruled out \\
Scale/lag robustness & 14/15 survivor observations robust & Survivor-class robustness is high & Robustness holds for all 19 observations \\
Diffuse tangential form & 13/14 diffuse all-tangential support & Diffuse tangential activity is the strongest repeated form & A universal edge trigger or signed force is identified \\
State moment closure & Median incremental $R^2$ near zero & The tested $(C,\dot{C},R)$ closure is not supported & Stochastic dynamics are impossible \\
Event-locality & Median near-pre event-minus-control effect about $-0.033$ & The tested state-matched near-pre route is not supported & Transitions have no special dynamics \\
Recent history & 14/19 above median shuffle, 6/19 above q95 & History effects are observation-specific & A universal memory rule is proven \\
Propagation & Confirmatory route not entered & Propagation remains open & No propagation exists \\
\bottomrule
\end{tabular}
\endgroup
\end{table*}
